%% ── Abstract ─────────────────────────────────────────────────
\chapter*{Abstract}
\addcontentsline{toc}{chapter}{Abstract}

The proliferation of sophisticated cyber-attacks has rendered static,
signature-based Intrusion Detection Systems (IDS) increasingly inadequate.
This thesis presents \caesar{} (\textit{Co-Evolutionary Adversarial Simulation
Engine for Attack and Response}), a novel AI-driven framework that
simultaneously simulates realistic cyber-attacks and trains adaptive defenses
through a closed co-evolutionary loop.

\caesar{} introduces four original algorithmic contributions.
First, the \textbf{Threat-Aware GAN (\tagan{})} conditions attack generation
on the current defense state, producing adversarial traffic that dynamically
adapts to deployed countermeasures — an approach absent from prior GAN-based
intrusion generation literature.
Second, the \textbf{Adaptive Defense Policy Network (\adpn{})} employs a
Dueling Double-DQN architecture to select countermeasures in real-time, driven
by a co-evolutionary reward signal that links defender performance to attacker
success.
Third, the \textbf{Temporal Attack Graph (\tagraph{})} accumulates a directed
knowledge graph of attack-defense interactions, enabling proactive defense
deployment with confidence-gated preemption — achieving statistically
significant threat anticipation.
Fourth, the \textbf{Manifold-Guided Adversarial Engine (\mgae{})} applies a
learned denoising diffusion process to project attack samples onto the benign
traffic manifold, producing adversarial examples that simultaneously evade
classifier-based and anomaly-based IDS.

Extensive experiments on the \textbf{real CICIDS2017} dataset (2.5M flows)
demonstrate that \caesar{} achieves a neutralization rate of
$\mathbf{91.9\% \pm 3.3\%}$, a robustness score of $\mathbf{0.967 \pm 0.007}$,
and a co-evolutionary gap of $\mathbf{+0.239}$, all statistically significant
at $\alpha = 0.05$ (Wilcoxon $p = 0.002$, Cohen's $d = 23.3$).
The \mgae{} engine achieves \textbf{complete evasion} of Random Forest and
Decision Tree classifiers on real data (detection rate: 98.5\% $\to$ 0\%).
The Autonomous Remediation Loop (\arl{}) demonstrates fully autonomous
incident response with \textbf{100\% healing success rate} over 300 simulation
ticks, maintaining a mean network health of \textbf{0.828}.

This research advances the state of the art in AI-driven cybersecurity by
providing the first unified framework for co-evolutionary attack simulation and
adaptive defense, with direct applicability to next-generation
\emph{self-healing} security systems.

\vspace{1em}
\noindent\textbf{Keywords:} Generative Adversarial Networks, Reinforcement Learning,
Intrusion Detection, Adversarial Machine Learning, Diffusion Models,
Co-evolutionary AI, Self-healing Systems, Cybersecurity.

\clearpage

%% ── Acknowledgements ─────────────────────────────────────────
\chapter*{Acknowledgements}
\addcontentsline{toc}{chapter}{Acknowledgements}

I would like to express my deepest gratitude to my supervisor
for their unwavering guidance, critical insight,
and support throughout this doctoral journey.

I am grateful to my colleagues at the University of Thessaly,
Department of Computer Science, for the stimulating
discussions and collaborative environment that shaped this research.

Finally, I dedicate this thesis to my family, whose encouragement made
everything possible.

\clearpage
